\documentclass[]{article}

\usepackage[italian]{babel}
\usepackage[margin=20mm, footskip = 20pt]{geometry}
\usepackage{array}
\usepackage{tabularx}
\usepackage{graphicx}
\usepackage{subfiles}
\usepackage{hyperref}
\usepackage{nameref}
\usepackage{titlesec}
\usepackage{longtable}
\usepackage[table]{xcolor}
\usepackage{titling}
\usepackage{lastpage}
\usepackage{ifthen}
\usepackage{calc}
\usepackage{soulutf8}
\usepackage{contour}
\usepackage{float}
\usepackage{fancyhdr}
\usepackage{multirow}
\usepackage{pgfgantt}
\usepackage{lscape}

\newcommand{\hr}{\par\vspace{-.1\ht\strutbox}\noindent\hrulefill\par}

\graphicspath{ {./}
	{./commons/res}
}

%--------------------------------------------------
% Comandi per inserire contenuto del documento
%--------------------------------------------------
\makeatletter

\newcommand\appendToGraphicsPath[1]{%
	\g@addto@macro\Ginput@path{{#1}}%
}

\newcommand{\setTitle}[1]{%
	\newcommand{\@phTitle}{#1}%
}
\newcommand{\phTitle}{\@phTitle}

\newcommand{\setDate}[1]{%
	\newcommand{\@phDate}{#1}%
}
\newcommand{\phDate}{\@phDate}

\newcommand{\setUso}[1]{%
	\newcommand{\@uso}{#1}%
}
\newcommand{\uso}{\@uso}

\newcommand{\setVersione}[1]{%
	\newcommand{\@versione}{#1}%
}
\newcommand{\versione}{\@versione}

\newcommand{\disabilitaVersione}{%
	\renewcommand{\setVersione}[1]{}%
	\renewcommand{\versione}{DISABILITATA}
}

\newcommand{\setResponsabile}[1]{%
	\newcommand{\@responsabile}{#1}%
}
\newcommand{\responsabile}{\@responsabile}

\newcommand{\setRedattori}[1]{%
	\newcommand{\@redattori}{#1}%
}
\newcommand{\redattori}{\@redattori}

\newcommand{\setVerificatori}[1]{%
	\newcommand{\@verificatori}{#1}%
}
\newcommand{\verificatori}{\@verificatori}

\newcommand{\setModifiche}[1]{%
	\newcommand{\@modifiche}{#1}%
}
\newcommand{\modifiche}{\@modifiche}

\makeatother 

%--------------------------------------------------
% Comandi per i documenti esterni e il glossario
%--------------------------------------------------

\newcommand{\dext}[1]{\textsc{#1\textsubscript{\textit{D}}}}

\newcommand{\glock}[1]{\textsc{#1\textsubscript{\textit{G}}}}

%--------------------------------------------------
% Comandi per impostare sottotitoli di quarto e quinto livello
%--------------------------------------------------

\setcounter{secnumdepth}{4}
\setcounter{tocdepth}{4}

\titleformat{\paragraph}
{\normalfont\normalsize\bfseries}{\theparagraph}{1em}{}
\titlespacing*{\paragraph}{0pt}{2.25ex plus 1ex minus .2ex}{1.5ex plus .2ex}

\titleformat{\subparagraph}
{\normalfont\normalsize\bfseries}{\thesubparagraph}{1em}{}
\titlespacing*{\subparagraph}{0pt}{1.75ex plus 1ex minus .2ex}{.75ex plus .1ex}

\appendToGraphicsPath{../../commons/res/}

%------------------------------
%
% COMANDI DI CONFIGURAZIONE
%
%------------------------------

\setTitle{Studio di Fattibilità}

\setVersione{0.2.0}

\setDate{14-12-2020}

\setResponsabile{Paolo Scanferlato}

\setRedattori{Alessandro Chimetto}

\setVerificatori{Giacomo Bulbarelli}

\setUso{Interno}

\setModifiche{
%%	Vers.	&	Nome				&	Ruolo		&	Data		&	Desrizione					\\%

%	M1		&	M2					&	M3			&	M3			&	M4							\\%
	0.2.0   &   Valton Tahiraj      &   Redattore	&	14-12-2020	&   Aggiunto C7\\
	0.1.0	&	Alessandro Chimetto	&	Redattore	&	12-12-2020	&	Stesura iniziale + studio C1}



\begin{document}

	% Direttive per la creazione del titolo tramite comando maketitle
\title{\huge \textsc{\phTitle{}} \\
	\vspace{11pt} \large \textsc{\phDate{}}}

\author{} % Non toccare
\date{} % Non toccare

%--------------------
% Frontespizio
%--------------------

% Logo del gruppo
\begin{figure}[t!]
	\centering
	\includegraphics[width=20em]{lclong}
\end{figure}

% Titolo / Nome
\maketitle
\thispagestyle{empty}

% Dati specifici sul doc in forma tabulare
\begin{table}[ht]
	\begin{center}
		\label{tab:Dati sul documento}
		\begin{tabular}{r|l}
			\multicolumn{2}{c}{ \textsc{Dati sul documento} } \\
			\hline
			\textbf{Versione} & \versione{} \\
			\textbf{Uso} & \uso{}  \\
			\textbf{Redattori} & \redattori{} \\
			\textbf{Verificatori} & \verificatori{} \\
			\textbf{Responsabile} & \responsabile{} \\
			\textbf{Destinatari} & lineCode \\
								& prof.\ Vardanega Tullio \\		
								& prof.\ Cardin Riccardo \\
			\ifthenelse{\equal{\uso}{Esterno}}{
								& Sanmarco Informatica
			}{} \\
		\end{tabular}
	\end{center}
\end{table}

\newpage

\renewcommand{\arraystretch}{2} % allarga le righe con dello spazio sotto e sopra
\begin{longtable}[H]{>{\centering\bfseries}m{2cm} >{\centering}m{3.5cm} >{\centering}m{2.5cm} >{\centering}m{3cm} >{\centering\arraybackslash}m{5cm}}
	\rowcolor{lightgray}
	{\textbf{Versione}} & {\textbf{Nominativo}} & {\textbf{Ruolo}} & {\textbf{Data}} & {\textbf{Descrizione}}  \\
	\endfirsthead%
	\rowcolor{lightgray}
	{\textbf{Versione}} & {\textbf{Nominativo}}  & {\textbf{Ruolo}} & {\textbf{Data}} & {\textbf{Descrizione}}  \\
	\endhead%
	\modifiche{}%
\end{longtable}

	\newpage

	%--------------------------------
	%
	% IL CONTENUTO INIZIA DA QUI
	%
	%--------------------------------

	%! INDICE

	\section{Introduzione}
		\subsection{Scopo del documento}
		Questo documento riporta lo studio svolto su ogni \glock{capitolato} proposto,
		enunciando per ognuno: finalità, tecnologie e caratteristiche ritenute
		negative o positive dai membri del gruppo.

		\subsection{Altri Documenti e Glossario}
		Terminologie specifiche all'interno del documento saranno scritte in \textsc{maiuscoletto} per evitare ambiguità e avranno a pedice:
		\begin{itemize}
			\item \textit{D} se indicano un documento specifico;
			\item \textit{G} se incluse nel \dext{glossario}.
		\end{itemize}

		\subsection{Riferimenti}
			\subsubsection{Normativi}
			\begin{itemize}
				\item \textbf{\dext{Norme di Progetto}}
			\end{itemize}

			\subsubsection{Informativi}
			\begin{itemize}
				\item \textbf{\glock{Capitolato} C1 - BlockCOVID: supporto digitale al contrasto della pandemia} \\
				(https://www.math.unipd.it/~tullio/IS-1/2020/Progetto/C1.pdf);

				\item \textbf{\glock{Capitolato} C2 - EmporioLambda: piattaforma di e-commerce in stile Serverless} \\
				(https://www.math.unipd.it/~tullio/IS-1/2020/Progetto/C2.pdf);

				\item \textbf{\glock{Capitolato} C3 - GDP: Gathering Detection Platform} \\ (https://www.math.unipd.it/~tullio/IS-1/2020/Progetto/C3.pdf);

				\item \textbf{\glock{Capitolato} C4 - HD Viz: visualizzazione di dati multidimensionali} \\
				(https://www.math.unipd.it/~tullio/IS-1/2020/Progetto/C4.pdf);

				\item \textbf{\glock{Capitolato} C5 - PORTACS: piattaforma di controllo mobilità autonoma} \\
				(https://www.math.unipd.it/~tullio/IS-1/2020/Progetto/C5.pdf);

				\item \textbf{\glock{Capitolato} C6 - RGP: Realtime Gaming Platform} \\
				(https://drive.google.com/file/d/1MQ8j4plXMsKjBfLPrUfHCv\_6pfnjcU5T/view?usp=sharing);

				\item \textbf{\glock{Capitolato} C7 - SSD: soluzioni di sincronizzazione desktop} \\
				(https://www.math.unipd.it/~tullio/IS-1/2020/Progetto/C7.pdf)
			\end{itemize}

	\newpage

	\section{Valutazione capitolato scelto}
		\subsection{C5 - PORTACS}
		%! AGGIUNGERE CAPITOLATO C5

	\newpage

	\section{Valutazione capitolati non scelti}
		\subsection{C1 - BlockCOVID}
			\subsubsection{Informazioni generali}
			\begin{itemize}
				\item \textbf{Titolo:} BlockCOVID;
				\item \textbf{Sottotitolo:} Supporto digitale al contrasto della pandemia;
				\item \textbf{Proponente:} Imola Informatica;
				\item \textbf{Committente:} Prof. Tullio Vardanega e Prof. Riccardo Cardin.
			\end{itemize}

			\subsubsection{Descrizione}
			Vista la corrente pandemia di SARS-CoV-2 e conseguenti accordi fra sindacati e imprese, si vuole creare un'applicazione in grado di tracciare la presenza del personale di un organizzazione (es. scuola o azienda) nelle proprie postazioni di lavoro e di tracciare l'igienizzazione di quest'ultime da parte dell'utente o un'azienda specializzata.

			\subsubsection{Finalità}
			Si vuole creare un'applicazione mobile che permetta all'utente di:
			\begin{itemize}
				\item segnalare, tramite tag \glock{RFID}, l'occupazione di una postazione di lavoro ad un server dedicato;
				\item prenotare una postazione libera ed igienizzata;
				\item segnalare se la postazione è stata pulita autonomamente alla fine dell'utilizzo.
			\end{itemize}
			L'azienda esterna specializzata deve poter utilizzare l'applicazione per:
			\begin{itemize}
				\item ricevere una lista delle postazioni da igienizzare;
				\item marcare la stanza come igienizzata.
			\end{itemize}
			Il server dedicato deve tenere traccia delle informazioni in una struttura dati di tipo \glock{blockchain} e deve mettere a disposizione dell'amministratore:
			\begin{itemize}
				\item un'interfaccia che consenta di creare, modificare o eliminare postazioni e stanze;
				\item la gestione del calendario delle prenotazioni;
				\item la possibilità di eseguire ricerche e creare report su personale addetto e igienizzazioni.
			\end{itemize}

			\subsubsection{Tecnologie}
			Parte dell'interesse di ricerca dell'azienda è l'esplorazione di nuove soluzioni tecnologiche dunque quasi tutte le tecnologie da essa proposte non sono obbligatorie ma solamente consigliate sulla base dell'esperienza:
			\begin{itemize}
				\item \textsc{\glock{API Rest} o \glock{gRPC}:} tecnologie di comunicazione asincrona fra applicazione e server;
				\item \textsc{\glock{Ethereum}:} tecnologia \glock{blockchain} per immagazzinare dati con opponibilità a terzi;
				\item \textsc{\glock{Kubernetes}, \glock{Openshift} o \glock{Rancher}:} orchestratori per gestire \glock{deploy} e scalabilità delle componenti server.
			\end{itemize}

			\subsubsection{Linguaggi di Programmazione}
			\begin{itemize}
				\item \textsc{\glock{Java}, \glock{Python} o \glock{Node.js}:} linguaggi molto diffusi, consigliati (ma non obbligatori) per lo sviluppo delle componenti \glock{back-end};
				\item \textsc{\glock{Kotlin} o \glock{Swift}:} linguaggi per lo sviluppo di applicazioni mobile rispettivamente su piattaforme \glock{Android} e \glock{iOS};
				\item \textsc{\glock{Solidity} o \glock{Vyper}:} linguaggi per lo sviluppo di \glock{smart contract} per la piattaforma \glock{Ethereum}.
			\end{itemize}

			\subsubsection{Vincoli di Progetto}
			\begin{itemize}
				\item Per la comunicazione asincrona fra applicazione e server, si richiede di utilizzare \glock{API Rest} o in alternativa \glock{gRPC};
				\item scansione dei codici nel tempo sufficiente a certificare la presenza della persona in postazione;
				\item presentazione di un resoconto su scelte e test effettuati per garantire una buona durata della batteria del dispositivo mobile nonostante l'utilizzo del lettore \glock{RFID};
			\end{itemize}

			\subsubsection{Aspetti positivi}
			\begin{itemize}
				\item Si sta affrontando un tema di grande attualità e la sua realizzazione, potenzialmente, implicherebbe un contributo anche a livello sociale;
				\item il \glock{capitolato} è fortemente politematico in quanto spazia dall'uso di hardware specializzato all'uso di linguaggi di programmazione attuali e strumenti rilevanti nell'ambito \glock{DevOps} per non parlare dell'utilizzo della struttura dati \glock{blockchain}; tutto ciò rende il lavoro proposto molto stimolante.
			\end{itemize}

			\subsubsection{Criticità}
			\begin{itemize}
				\item Il fatto che questo \glock{capitolato} sia estremamente politematico implica che la mole di lavoro nel progetto da esso derivato sia imponente. Ci sono molte tecnologie nuove da imparare in poco tempo e al fine di soddisfare i soli requisiti obbligatori si richiede lo sviluppo di due applicazioni di più che modeste dimensioni (mobile e server) che implementino tutte le tecnologie;
				\item non è chiaro se l'hardware specifico per tag \glock{RFID} venga fornito dall'azienda e in che quantità; ciò potrebbe diventare un problema nel momento in cui le unità a disposizione fossero poche in quanto poche persone ne avrebbero accesso e, vista la mancanza di contatti causa pandemia, non ci sarebbe la possibilità di scambiarle con altri membri del gruppo.
			\end{itemize}

			\subsubsection{Conclusioni}
			Nonostante il questo \glock{capitolato} abbia suscitato un forte interesse in tutti i membri del gruppo per questioni sia tecnologiche che sociali, si teme che la mole di lavoro sia troppo ampia e di conseguenza che non si riesca a consegnare il prodotto in tempi accettabili per il gruppo.

		\newpage

		\subsection{C2 - Emporio Lambda}
		
		\subsubsection{Informazioni generali}
			\begin{itemize}
				\item \textbf{Titolo:} Emporio Lambda;
				\item \textbf{Proponente:} Red Babel;
				\item \textbf{Committente:} Prof. Tullio Vardanega e Prof. Riccardo Cardin.
			\end{itemize}

			\subsubsection{Descrizione}
			Le piattaforme di e-commerce possono essere triviali, come pure complesse, in base alle funzionalità che si vogliono garantire. Inoltre i bisogni delle aziende possono variare ampiamente nel tempo e si scopre quindi la necessità di appoggiarsi a tecnologie affidabili per l'uso richiesto, e che garantiscano adattabilità a cambiamenti futuri.

			\subsubsection{Finalità}
			Emporio Lambda mira ad essere una piattaforma di e-commerce con architettura serverless appoggiandosi ad \glock{AWS}. Si vogliono garantire le funzionalità tipiche di una piattaforma commerciale, che permetta agli amministratori un'adeguata gestione degli ordini e dell'inventario, ed agli utenti un'adeguata esperienza di ricerca ed acquisto dei prodotti disponibili. 
			\\
			Il sito dovrà essere performante e garantire agli utenti un accesso immediato alla ricerca dei prodotti. Sarà importante l'uso di tecniche di \glock{SEO} soprattutto per le pagine dei prodotti in vendita. 
			\\
			Dovrà essere disponibile uno strumento di autenticazione, che potrà essere realizzato appoggiandosi ad \glock{Amazon Cognito}.

			\subsubsection{Tecnologie}
			\begin{itemize}
				\item \textsc{AWS Lambda:} piattaforma di calcolo serverless, fornita come parte di \glock{AWS} che si occupa della gestione \glock{back-end} di applicazioni web;
				\item \textsc{Amazon CloudWatch:} servizio di monitoraggio e osservabilità che raccoglie dati di monitoraggio e operativi sotto forma di log, parametri ed eventi, fornendo una visualizzazione unificata delle risorse;
				\item \textsc{Amazon Cognito:} servizio la gestione di registrazione e controllo degli accessi a piattaforme web.
			\end{itemize}

			\subsubsection{Linguaggi di Programmazione}
			\begin{itemize}
				\item \textsc{\glock{Next.js}:} framework per lo sviluppo \glock{front-end} di siti web;
				\item \textsc{\glock{Typescript}:} linguaggio che va ad estendere \glock{Javascript} dotandolo di tipizzazione;
				\item \textsc{\glock{Node.js}:} linguaggio molto diffuso, per lo sviluppo delle componenti \glock{back-end} da appoggiare ad \glock{AWS}.
			\end{itemize}
			

			\subsubsection{Vincoli di Progetto}
			Il sito deve permettere il supporto ai ruoli di: 
			\begin{itemize}
			    \item amministratore di sistema: per la configurazione ed  amministrazione del sito;
			    \item venditore: per l'organizzazione e la gestione dei prodotti e delle loro informazioni;
			    \item cliente: utente generico del sito che può visualizzare ed acquistare i prodotti.
			\end{itemize}				
		A livello di funzionalità per il cliente, devono essere previste le seguenti pagine:  homepage, profilo utente, elenco prodotti, dettagli prodotti, carrello, checkout.\\
		Si deve poi garantire un'adeguata integrazione con i servizi amministrativi tipici di questo ambiente, come: supporto utenti, tracciamento delle operazioni, inventario magazzino.\\

			\subsubsection{Aspetti positivi}
			\begin{itemize}
				\item le tecnologie richieste per la realizzazione del \glock{capitolato} sono state giudicate di interesse;
				\item la difficoltà generale dei requisiti è considerata adeguata per i componenti del gruppo.
			\end{itemize}

			\subsubsection{Criticità}
			\begin{itemize}
			    \item l'e-commerce è stato considerato come un target non stimolante da parte del gruppo;
				\item il materiale reso disponibile dal proponente è stato giudicato non molto chiaro, e si ha il timore che ci potrebbero essere problemi di cooperazione in futuro.
			\end{itemize}

			\subsubsection{Conclusioni}
			Nonostante le tecnologie che si andrebbero ad usare sembrino interessanti, nel complesso il \glock{capitolato} non ha suscitato interesse sufficiente rispetto alle altre opzioni disponibili al gruppo.
		
		\newpage
		%! AGGIUNGERE RESTANTI CAPITOLATI E CONCLUSIONE FINALE
		
		\subsection{C3 - GDP}
		\subsubsection{Informazioni generali}
		\begin{itemize}
			\item \textbf{Titolo:} GDP;
			\item \textbf{Sottotitolo:} Gathering Detection Platform;
			\item \textbf{Proponente:} Sync Lab;
			\item \textbf{Committente:} Prof. Tullio Vardanega e Prof. Riccardo Cardin.
		\end{itemize}
		
		\subsubsection{Descrizione}
		In seguito alla pandemia del virus COVID-19 scoppiata negli ultimi mesi sono sorti numerosi problemi. Dopo un periodo di quarantena forzata la vita doveva ripartire ed ai cittadini è stata permessa la circolazione. Durante la vita quotidiana si generano spesso assembramenti in numerosi luoghi e per non aggravare la situazione bisogna cercare di prevenire questo problema. Grazie all’informatica possiamo quindi scongiurare l’aggravarsi della situazione dovuta ai problemi sociali legati all’affollamento.
		
		\subsubsection{Finalità}
		
		Si vuole creare un’applicazione web che permetta di:
		\begin{itemize}		
		\item misurare tramite dei contapersone e date le immagini/stream delle videocamere di contare le persone a bordo dei mezzi pubblici;
		\item poter acquisire continuativamente nel tempo e in modalità a bassa latenza delle informazioni raccolte da dispositivi di flussi dati;	
		\item elaborare in tempo reale i dati acquisiti per rappresentare le variazioni nel tempo dei dati monitorati, come ad esempio la variazione dei flussi delle persone in ambienti urbani. Confrontare e correlare tra loro dati provenienti da flussi diversi, ad esempio correlando i dati raccolti a bordo dei mezzi di trasporto di interesse con i dati raccolti in prossimità delle aree di interesse quali ad esempio fermate di autobus, strade dedicate al traffico veicolare, zone urbane. E' poi fondamentale archiviare tutti i dati acquisiti, ed i risultati delle loro elaborazioni affinché siano conservati nel tempo e comunque disponibili in qualunque caso sia richiesto il loro utilizzo; 
		\item \textsc identificazione, a partire dai dati acquisiti ed elaborati, di eventi che nel tempo sono risultati aver concorso ad esempio all’insorgere di alterazioni significative del flusso di utenti;
		\item \textsc previsione dell’insorgenza futura di variazioni significative di flussi di persone al fine di poter fornire degli indicatori automatici in grado di fornire supporto alle decisioni da parte del personale tecnico specializzato preposto a tale scopo.
		\end{itemize}
		 	
		\subsubsection{Tecnologie}
		Ci sono alcune scelte preferenziali per lo svolgimento del progetto:
		\begin{itemize}		
		\item utilizzo di \glock{Java} e \glock{Angular} per lo sviluppo delle parti di \glock{Back-end} e \glock{Front-end} della componente Application del sistema;
		\item il framework \glock{Leaflet} (https://leafletjs.com) per la gestione delle mappe (heatmap);
		\item utilizzo di protocolli asincroni per le comunicazioni tra diverse componenti;
		\item utilizzo del pattern Publisher/Subscriber e adozione del protocollo \glock{MQTT}(‘MQ Telemetry Transport or Message Queue Telemetry Transport’).
		\end{itemize}
		
		\subsubsection{Aspetti positivi}
		Il progetto risultava secondo il gruppo uno dei tre capitolati più interessanti e stimolanti:
		\begin{itemize}			
		\item il problema da affrontare è qualcosa che riguarda la vita di tutti i giorni e che potrebbe portare a miglioramenti anche fuori da una situazione di emergenza COVID-19;
		\item le tecnologie richieste per risolvere il problema e la componente di IA rappresenta un’ottima occasione per ampliare il bagaglio di conoscenze.	
		\end{itemize}
		
		\subsubsection{Criticità}
		\begin{itemize}
		\item La parte di IA e la problematica da risolvere sembrava meno interessante rispetto al capitolato scelto.
		\end{itemize}
		
		\subsubsection{Conclusioni}
		Dopo aver preso ampiamente in considerazione il capitolato in questione ed aver seguito la presentazione ed il seminario tecnologico il gruppo ha preferito orientarsi verso un’altra scelta che sembrava più stimolante. 
		
		\newpage 
		
		
		\subsection{C7 - SSD}
		\subsubsection{Informazioni generali}
		\begin{itemize}
			\item \textbf{Titolo:} SSD;
			\item \textbf{Sottotitolo:} soluzioni di sincronizzazione desktop;
			\item \textbf{Proponente:} ZEXTRAS;
			\item \textbf{Committente:} Prof. Tullio Vardanega e Prof. Riccardo Cardin.
		\end{itemize}
		
		\subsubsection{Descrizione}
		Si vuole approfondire il mondo degli algoritmi di sincronizzazione desktop, essendo la piattaforma principale per molti utenti professionali, in un modo non triviale che non si limiti a permettere il lavoro in un dispositivo mantenendo una copia del lavoro salvata su \glock{Cloud}.
		
		\subsubsection{Finalità}
		I principali obbiettivi di questo progetto sono:
		\begin{itemize}
			\item Sviluppo di un algoritmo solido ed efficiente in grado di garantire il salvataggio in \glock{Cloud} del lavoro e contemporaneamente la sincronizzazione dei cambiamenti presenti in   \glock{Cloud};
			\item Sviluppo di un'interfaccia multipiattaforma per l'uso dell'algoritmo nei più importanti sistemi operativi desktop esistenti(\glock{MacOs},\glock{Windows},\glock{Linux})
			\item Utilizzo dell'algoritmo sviluppato per richiedere e fornire i cambiamenti ai contenuti in sincronizzazione verso il prodotto del proponente \glock{Zextras Drive}.
		\end{itemize}
		Le principali funzioni richieste sono:
		\begin{itemize}
			\item Configurazione ed autenticazione dell'utente;
			\item Gestione di cosa sincronizzare e di cosa ignorare nelle cartelle \glock{Cloud};
			\item Gestione di cosa sincronizzare e di cosa ignorare nelle cartelle locali;
			\item Sincronizzazione costante dei cambiamenti, siano essi locali o remoti;
			\item Possibilità di modifica delle preferenze a posteriori;
			\item Sistema di notifica utente dei cambiamenti;
		\end{itemize}
		Con l'aggiunta di funzionalità avanzate quali:
		\begin{itemize}
			\item Gestione delle condivisioni;
			\item Integrazione protocollo \glock{MAPI}; 
			\item Integrazione con il prodotto Web;
		\end{itemize} 	
		\subsubsection{Tecnologie}
		Il proponente sancisce che la soluzione sviluppata non dipenda da installazione di \glock{Framework} terzi per funzionare.
		Consiglia quindi l'uso dei seguenti \glock{Framework}: 
		\begin{itemize}
			\item \textsc{\glock{Qt Framework}}: \glock{Framework} consigliato per lo sviluppo dell'interfaccia e del controller d'architettura;
			\item \textsc{\glock{Python}}: Linguaggio consigliato per lo sviluppo della Backend Business Logic;
			
		\end{itemize}
		
		
		\subsubsection{Aspetti positivi}
		\begin{itemize}
			\item Le tecnologie sono conosciute dalla maggior parte dei membri del gruppo.
		\end{itemize}
		
		\subsubsection{Criticità}
		\begin{itemize}
			\item La documentazione del \glock{capitolato} risulta scritta in modo approssimato facendo perdere al gruppo l'entusiasmo di lavorare con questo proponente;
			\item L'idea risulta poco innovativa.
		\end{itemize}
		
		\subsubsection{Conclusioni}
		Questo \glock{capitolato} non ha suscitato interesse in nessun membro del gruppo. Il gruppo ha quindi deciso di  concentrarsi su altri \glock{capitolati}. 
		
		\newpage
\end{document}

